%COMANDO CHE DETERMINA IL TIPO DI DOCUMENTO CHE SI VUOLE CREARE

\documentclass[a4paper,12pt]{report}

%ELENCO DEI PACCHETTI UTILI PER LA STESURA DEL DOCUMENTO

\usepackage[utf8]{inputenc}
\usepackage[english, italian]{babel}
\usepackage{graphicx}
\usepackage{float}
\usepackage{tabularx}
\usepackage{makecell}
\usepackage{titlesec}
\usepackage{fancyhdr}
\usepackage{lastpage}
\usepackage{xurl}
\usepackage{hyperref}
\usepackage{geometry}
\usepackage[table,dvipsnames]{xcolor}
\setcounter{tocdepth}{5}
\setcounter{secnumdepth}{5}
\usepackage{caption}
\usepackage{etoolbox}% >= v2.1 2011-01-03
\usepackage{tikz}
\usepackage{listings}
\usepackage{scalerel}
\usepackage{longtable}
\usepackage{comment}
\usepackage{amssymb}
\usepackage{listings}

\titleformat{\chapter}[display]
{\normalfont\bfseries}{}{0pt}{\LARGE}

%COMANDO PER LA SPAZIATURA DEI TITOLI DAL BORDO DEL FOGLIO

\titlespacing*{\chapter}{0cm}{0cm}{0.2cm}

%COMANDO PER LA SPAZIATURA DEL TESTO DAI BORDI LATERALI

\geometry{
	left=20mm,
	right=20mm,
}

%COMADNO PER AVERE L'INDICE DEL NOME CHE SI VUOLE

\renewcommand{\contentsname}{Indice}

%COMADNI PER OTTENERE SUBSUBSECTION NUMERATE E PRESENTI NELL'INDICE

\setcounter{tocdepth}{5}
\setcounter{secnumdepth}{5}

%COMANDI PER OTTENERE HEADER E FOOTER

\pagestyle{plain}

\fancypagestyle{plain}{
	\fancyhf{}
	\lhead{\includegraphics[width=3cm]{../../immagini/minilogo.jpg}}
	\chead{}
	\rhead{\fontsize{12}{10}Verbale Esterno 2021-04-15}
	\lfoot{}
	\cfoot{\thepage\ di \pageref{LastPage}}
	\rfoot{}
}

\pagestyle{plain}

%COMANDI PER LINK

\hypersetup{
	colorlinks=true,
	linkcolor=black,
	filecolor=black,
	urlcolor=blue,
	citecolor=black,
}

%COMANDI PER SNIPPET DI CODICE

\BeforeBeginEnvironment{lstlisting}{\begin{mdframed}\vspace{-0.7em}}
	\AfterEndEnvironment{lstlisting}{\vspace{-0.5em}\end{mdframed}}

% needed for \lstcapt
\def\ifempty#1{\def\temparg{#1}\ifx\temparg\empty}

% make new caption command for listings
\usepackage{caption}
\newcommand{\lstcapt}[2][]{%
	\ifempty{#1}%
	\captionof{lstlisting}{#2}%
	\else%
	\captionof{lstlisting}[#1]{#2}%
	\fi%
	\vspace{0.75\baselineskip}%
}

\definecolor{atomlightorange}{rgb}{0.88,0.76,0.55}
\definecolor{atomdarkgrey}{RGB}{59,62,75}

% set listings
\lstset{%
	basicstyle=\footnotesize\ttfamily\color{atomlightorange},
	framesep=20pt,
	belowskip=10pt,
	aboveskip=10pt
}

% add frame environment
\usepackage[%
framemethod=tikz,
skipbelow=8pt,
skipabove=13pt
]{mdframed}
\mdfsetup{%
	leftmargin=0pt,
	rightmargin=0pt,
	backgroundcolor=atomdarkgrey,
	middlelinecolor=atomdarkgrey,
	roundcorner=6
}

%BIBLIOGRAFIA

\makeatletter
\def\thebibliography#1{\chapter*{Bibliografia\@mkboth
		{Bibliografia}{Bibliografia}}\list
	{[\arabic{enumi}]}{\settowidth\labelwidth{[#1]}\leftmargin\labelwidth
		\advance\leftmargin\labelsep
		\usecounter{enumi}}
	\def\newblock{\hskip .11em plus .33em minus .07em}
	\sloppy\clubpenalty4000\widowpenalty4000
	\sfcode`\.=1000\relax}
\makeatother

%DEFINIZIONE COLORI

\definecolor{airforceblue}{rgb}{0.36, 0.54, 0.66}


\begin{document}

\makeatletter
\begin{titlepage}
	\begin{center}
		\vspace*{-4,0cm}
		\author{Jawa Druids} 
		\title{Nome documento}
		\date{} %LASCIARE QUESTO CAMPO VUOTO, SE LO TOLGO STAMPA LA DATA CORRENTE
		\includegraphics[width=0.7\linewidth]{../immagini/DRUIDSLOGO.jpg}\\[4ex]
		{\huge \bfseries  \@title }\\[2ex] 
		{\LARGE  \@author}\\[50ex]
		\vspace*{-8,5cm}
		\begin{table}[H]
			\renewcommand{\arraystretch}{1.4}
			\centering
			\begin{tabular}{r | l}
				\textbf{Versione} & x.x.x \\%RIGA PER INSERIRE LA VERSIONE ULTIMA DEL DOCUMENTO
				\textbf{Data approvazione} & xx-xx-xxxx\\
				\textbf{Responsabile} & Nome Cognome\\
				\textbf{Redattori} & Nome Cognome \\
				\textbf{Verificatori} & \makecell[tl]{Nome Cognome \\ Nome Cognome} \\
				%MAKECELL SERVE PER POI ANDARE A CAPO ALL'INTERNO DELLA CELLA
				\textbf{Stato} & Stato\\
				\textbf{Lista distribuzione} & \makecell[tl]{Nome Gruppo \\ Nome Professori \\ Nome Azienda Proponente}\\
				\textbf{Uso} & Uso del documento            
			\end{tabular}
		\end{table}
		\vspace{0.1cm}
		\hfill \break
		\fontsize{17}{10}\textbf{Sommario} \\
		\vspace{0.1cm}
		Brevissima descrizione del documento 
	\end{center}
\end{titlepage}
\makeatother

\quad
\begin{center}
	\LARGE\textbf{Registro delle modifiche}
\end{center}

\def\tabularxcolumn#1{m{#1}}
{\rowcolors{2}{RawSienna!90!RawSienna!20}{RawSienna!70!RawSienna!40}

\begin{center}
	\renewcommand{\arraystretch}{1.4}
\begin{tabularx}{\textwidth}{|X|c|c|c|c|}
	\hline
	\rowcolor{airforceblue}
	\textbf{Modifica} & \textbf{Autore} & \textbf{Ruolo} & \textbf{Data} & \textbf{Versione}\\
	\hline
	\textit{Descrizione delle modifiche} & Nome Cognome & \textit{Analista} & 23-11-2020 & v0.0.3 \\
	\hline
	\textit{Descrizione delle modifiche} & Nome Cognome & \textit{Analista} & 20-11-2020 & v0.0.2 \\
	\hline
	\textit{ Stesura del documento} & Emma Roveroni & \textit{Analista} & 06-12-2020 & v0.0.1 \\
	\hline
\end{tabularx}
\end{center}

%COMANDO PER LA CREAZIONE DELL'INDICE

\tableofcontents{}

%PER RENDERE PIÙ CHIARA LA STESURA DEI DOCUMENTI È MEGLIO LASCIARE SEPARATI IN FILE DIVERSI OGNI CAPITOLO

% \input{esempio} -- esempio di codice per inserire un nuovo capitolo

\end{document}